% Exported from Math++. Compile with pdflatex.
\documentclass[11pt,a4paper]{article}
\usepackage[T1]{fontenc}
\usepackage[utf8]{inputenc}
\usepackage{lmodern}
\usepackage{amsmath,amssymb,amsthm,mathtools}
\usepackage{graphicx,booktabs,xcolor}
\usepackage[hidelinks]{hyperref}
\usepackage[margin=25mm]{geometry}
\theoremstyle{plain}
\newtheorem{theorem}{Theorem}
\newtheorem{lemma}[theorem]{Lemma}
\newtheorem{corollary}[theorem]{Corollary}
\newtheorem{proposition}[theorem]{Proposition}
\newtheorem{conjecture}[theorem]{Conjecture}
\theoremstyle{definition}
\newtheorem{definition}[theorem]{Definition}
\newtheorem{example}[theorem]{Example}
\theoremstyle{remark}
\newtheorem*{remark}{Remark}
\setlength{\parindent}{0pt}\setlength{\parskip}{0.6em}
\begin{document}
\section*{Showing numbers your way}

Results are worked out exactly. The \texttt{\%} lines below only change how they are shown, never what they are.

A fraction stays a fraction: $\frac{22}{7} + 0 = \frac{22}{7}$

The same answer as a decimal with four significant figures:

\begin{equation*}
  \frac{22}{7} + 0 = 3.143
\end{equation*}
\begin{equation*}
  \sqrt{2} = 1.414
\end{equation*}
Two decimal places for the next result only:

\begin{equation*}
  \pi \approx 3.14
\end{equation*}
\begin{equation*}
  \pi \approx 3.142
\end{equation*}
Back to exact, with large numbers in scientific notation, and then engineering notation:

\begin{equation*}
  6.02 \times 10^{23} + 0 = 6.02 \times 10^{23}
\end{equation*}
\begin{equation*}
  123456789 \cdot 1000 = 123.456789 \times 10^{9}
\end{equation*}
Mixed fractions:

\begin{equation*}
  \frac{7}{2} + 0 = 3\frac{1}{2}
\end{equation*}
Angles from inverse trigonometric functions, in degrees:

\begin{equation*}
  \arcsin(1/2) = 30^\circ
\end{equation*}
Complex numbers in polar form:

\begin{equation*}
  (3 + 4i) + 0 = 5\angle 53.13010235^\circ
\end{equation*}
\end{document}
